\documentclass[]{article}
\usepackage[T1]{fontenc} 
\usepackage[icelandic]{babel}

\begin{document}
Fyrsta verkefnið sem mér var falið í upphafi tímabils var að breyta skipulagi orðaskrárnar. Orðaskrá íslenska stærðfræðifélagsins var upprunalega sett upp sem ensk-íslensk orðabók en ég var beðinn um að breyta henni þannig að hún sé í staðinn íslensk íðorðabók. Til þess fékk ég aðgang að gagnaskránni \texttt{ord8.dat} sem inniheldur orðaskránna á óstöðluðu sniði í formi \LaTeX skipana. Henni fylgdi engin skjölun og í augnablikinu virðist ekki vera til neinn þáttari fyrir \LaTeX kóða ætlaður til almennrar notkunar. Þess vegna ákvað ég að skrifa eigin þáttara fyrir snið orðaskrárnar frá grunni sem tók um það bil eina viku þar til hann var orðinn nothæfur (þann kóða má finna í \texttt{parser} möppunni). Það hefði ábyggilega tekið svipað langan tíma að læra að nota tilbúinn \LaTeX þáttara þó svo að hann hefði verið til og að vera með eigin þáttara hefur marga kosti. 

Þegar þáttarinn var tilbúinn varð mjög auðvelt að ``snúa við'' orðaskránni þannig að íslensk íðorð væru í forgrunni (kóðinn sem gerir það má finna í \texttt{invert.py}. Fyrsta tilraunin til þess var hins vegar ekki mjög árangursrík sökum þess að í mörgum tilfellum voru ensk samheiti þýdd eins í íslensku. Það olli því að til urðu margar færslur sem náðu utan um sama íðorðið. Því þurfti að skrifa kóða til þess leita að og eyða svona tvítekningum (sá kóða má finna í \texttt{remove\_duplicates.py}). Erfitt er að fanga nákvæmlega hvað það þýðir að tvær færslur séu tvítekningar og því erfitt að leysa þetta vandamál algjörlega sjálfvirkt. Fyrsta tilraunin reyndist vera of aggressíf og eyddi út sumum færslum sem voru ekki tvítekningar. Þegar það kom í ljós var reikniritið gert örlítið fágaðra en þá varð það örlítið of íhaldssamt. Tekin var ákvörðun að það væri skárra en að íðorðum væri hent út og ákveðið að hætta að vinna í reikniritinu. Því eru enn tvítekningar í orðaskránni sem þarf að fara yfir handvirkt. 

Samhliða þessu var skrifað kóða sem prentar út orðaskránna á CSV formi (sem má finna í \texttt{parser/ast.py}). Eftir að orðaskráin var komin á það form var notað Greyni (málgreinir fyrir íslensku frá Miðeind) til þess að greina kyn og orðaflokk íðorðana sjálfvirkt. Það reyndist ófullkomið og því þurfti að fara handvirkt yfir og laga villur á allri orðaskránni sem tók ábyggilega 1-1.5 viku. Þegar það var búið var útbúið kóða til að prenta út orðaskrána á formi \LaTeX töflu (sá kóða má finna í \texttt{generate\_tex\_table.py}, og töfluna sjálfa má finna í \texttt{ordaskra\_table.tex}). Hjálp gervigreindar var fengin við að skrifa kóðann sem býr til töfluna. Það er vegna þess að inntakið í þann kóða er CSV skráin en ekki \texttt{ord8.dat}. Það var litið á það sem svo að það væri of mikil tímaeyðsla að læra almennilega á \texttt{Pandas} skráasafnið enda var þessi tafla einungis ætluð til að einfalda vinnuna við það að skrifa skilgreiningar. Gervigreindin var líka fengin til að skrifa megnið af kóðanum í \texttt{process\_table.py} sem tekur þessa \LaTeX töflu og síar út þær færslur sem er búið að skilgreina til þess að hægt sé að hafa yfirlit yfir stöðu verkefnisins (það framleiðir nýja töflu sem má finna í \texttt{ordaskra\_table\_finished.text}. Því miður stóð gervigreindin sig ekki nógu vel og þessum kóða er mjög ábótavant. Eins og áður var getið eru þessar töflur hins vegar bara til að hjálpa við vinnuna og því ekki mikilvægt að þessi kóði sé réttur.

Í framhaldi var snúið athyglinni algjörlega að skilgreiningarvinnuni og fór meirihluti tímabilsins í hana. Ákveðið var að einblína á efni sem er kennt í framhaldsskólum og þá sérstaklega rúmfræði. Þegar vinnan var komin vel á stað var síðan ákveðið að taka skilgreiningar sem voru nú þegar til á stae.is og aðlaga þær að orðsaskránni en í lok tímabils var búið að koma meira og minna öllum þeim skilgreiningum inn í orðaskránna. Því næst var skoðað STÆ 103 kennslubókina eftir Jón Hafstein Jónsson, Níels Karlsson og Stefán G. Jónsson og skilgreint flestöll hugtökin sem fjallað var um í þeirri bók. Í heildina var skilgreint um það bil þrjú hundruð íðorð.

Í lok tímabils var síðan séð til þess að orðaskráin væri á birtingarhæfu formi og skrifað var kóða til þess að færa skilgreiningarnar úr \LaTeX töflunni aftur í ``vinnsluútgáfuna'' af CSV skránni (sá kóða má finna í \texttt{process\_table.py}). Það sem einkennir þessa vinnsluútgáfu helst er að íðorðið og öll samheiti þess er búið að kremja í einn dálk. Bætt var við nýjum dálki í skránna sem er ætlaður til að hægt sé að velja hvert af þessum orðum eigi að standa fyrir íðorðið. Þegar CSV skráin var fyrst búin til var valið fyrstu íslensku þýðinguna sem íðorðið en í mörgum tilfellum er það ekki algengasta orðið sem er notað yfir hugtakið sem um ræðir. Skrifaður var kóði sem tekur tillit til þessa dálks og velur tilgreinda orðið sem íðorðið (sá kóða má finna í \texttt{select\_word.py}).  
\end{document}